\documentclass[10pt,a4paper]{amsart}
\usepackage{amsmath,amssymb,amsthm,graphicx}
\usepackage{verbatim}
\newenvironment{reference}{\comment}{\endcomment}
%\newenvironment{reference}{}{}
\newenvironment{slogan}{\comment}{\endcomment}
\newenvironment{history}{\comment}{\endcomment}

% For commutative diagrams we use Xy-pic
\usepackage[all]{xy}

% We use 2cell for 2-commutative diagrams.
\xyoption{2cell}
% Theorem environments.
%
\theoremstyle{plain}
\newtheorem{theorem}[subsection]{Theorem}
\newtheorem{proposition}[subsection]{Proposition}
\newtheorem{lemma}[subsection]{Lemma}

\theoremstyle{definition}
\newtheorem{definition}[subsection]{Definition}
\newtheorem{example}[subsection]{Example}
\newtheorem{exercise}[subsection]{Exercise}
\newtheorem{situation}[subsection]{Situation}

\theoremstyle{remark}
\newtheorem{remark}[subsection]{Remark}
\newtheorem{remarks}[subsection]{Remarks}

\numberwithin{equation}{subsection}

% Macros
%
\def\lim{\mathop{\mathrm{lim}}\nolimits}
\def\colim{\mathop{\mathrm{colim}}\nolimits}
\def\Spec{\mathop{\mathrm{Spec}}}
\def\Hom{\mathop{\mathrm{Hom}}\nolimits}
\def\Ext{\mathop{\mathrm{Ext}}\nolimits}
\def\SheafHom{\mathop{\mathcal{H}\!\mathit{om}}\nolimits}
\def\SheafExt{\mathop{\mathcal{E}\!\mathit{xt}}\nolimits}
\def\Sch{\mathit{Sch}}
\def\Mor{\mathop{\mathrm{Mor}}\nolimits}
\def\Ob{\mathop{\mathrm{Ob}}\nolimits}
\def\Sh{\mathop{\mathit{Sh}}\nolimits}
\def\NL{\mathop{N\!L}\nolimits}
\def\CH{\mathop{\mathrm{CH}}\nolimits}
\def\proetale{{pro\text{-}\acute{e}tale}}
\def\etale{{\acute{e}tale}}
\def\QCoh{\mathit{QCoh}}
\def\Ker{\mathop{\mathrm{Ker}}}
\def\Im{\mathop{\mathrm{Im}}}
\def\Coker{\mathop{\mathrm{Coker}}}
\def\Coim{\mathop{\mathrm{Coim}}}

% Boxtimes
%
\DeclareMathSymbol{\boxtimes}{\mathbin}{AMSa}{"02}

%
% Macros for moduli stacks/spaces
%
\def\QCohstack{\mathcal{QC}\!\mathit{oh}}
\def\Cohstack{\mathcal{C}\!\mathit{oh}}
\def\Spacesstack{\mathcal{S}\!\mathit{paces}}
\def\Quotfunctor{\mathrm{Quot}}
\def\Hilbfunctor{\mathrm{Hilb}}
\def\Curvesstack{\mathcal{C}\!\mathit{urves}}
\def\Polarizedstack{\mathcal{P}\!\mathit{olarized}}
\def\Complexesstack{\mathcal{C}\!\mathit{omplexes}}
% \Pic is the operator that assigns to X its picard group, usage \Pic(X)
% \Picardstack_{X/B} denotes the Picard stack of X over B
% \Picardfunctor_{X/B} denotes the Picard functor of X over B
\def\Pic{\mathop{\mathrm{Pic}}\nolimits}
\def\Picardstack{\mathcal{P}\!\mathit{ic}}
\def\Picardfunctor{\mathrm{Pic}}
\def\Deformationcategory{\mathcal{D}\!\mathit{ef}}
\usepackage{cleveref}
\crefname{originaldefinition}{original definition}{original definitions}
\Crefname{originaldefinition}{Original definition}{Original definitions}
\makeatletter
\newlabel{definition-cycles}{{02QR}{source}{}{sourcepointer-1}{}}
\newlabel{definition-cycles@cref}{{[originaldefinition][0][]02QR}{[1][0][]source}{}{}{}}
\makeatother
\crefname{originallemma}{original lemma}{original lemmas}
\Crefname{originallemma}{Original lemma}{Original lemmas}
\makeatletter
\newlabel{lemma-push-proper-bivariant}{{0EPK}{source}{}{sourcepointer-2}{}}
\newlabel{lemma-push-proper-bivariant@cref}{{[originallemma][0][]0EPK}{[1][0][]source}{}{}{}}
\makeatother
\crefname{originallemma}{original lemma}{original lemmas}
\Crefname{originallemma}{Original lemma}{Original lemmas}
\makeatletter
\newlabel{lemma-exterior-product-well-defined-dim-1}{{0FC4}{source}{}{sourcepointer-3}{}}
\newlabel{lemma-exterior-product-well-defined-dim-1@cref}{{[originallemma][0][]0FC4}{[1][0][]source}{}{}{}}
\makeatother
\crefname{originallemma}{original lemma}{original lemmas}
\Crefname{originallemma}{Original lemma}{Original lemmas}
\makeatletter
\newlabel{lemma-bivariant-zero}{{02UC}{source}{}{sourcepointer-4}{}}
\newlabel{lemma-bivariant-zero@cref}{{[originallemma][0][]02UC}{[1][0][]source}{}{}{}}
\makeatother
\title{Operational Chow cohomology of a scheme over a base of dimension one is Chow homology in complementary degrees}
\author{The Stacks Project Authors}
\date{stacks-project:chow.tex:local-snapshot-20261012; GNU FDL 1.2 or later}
\begin{document}
\maketitle
\noindent\textit{Editorial note.} All delivered mathematical text is the author's own, copied from the pinned Stacks Project chapter file chow.tex; no statement or proof step was rewritten, reordered or removed.\par\medskip
\noindent\textit{Editorial note.} Preamble provenance: the pinned chapter file opens with an input of the shared preamble and carries no preamble text of its own, so the project's single Stacks-wide preamble staged at sources/stacks-chow/preamble.tex is excerpted verbatim as two spans, lines 14-18 (Xy-pic setup) and lines 154-217 (the native theorem declarations and mathematical macros). Nothing was added, removed or redefined; the wrapper class and the amsmath/amssymb/amsthm/graphicx base packages are supplied by the assembler.\par\medskip
\noindent\textit{Editorial note.} The counted statement and proof contain no cross reference, no citation and no included file, so the delivered body is self contained.\par\medskip
\noindent\textit{Editorial note.} Canonical location of the counted proposition: Stacks tag 0FC5 (sources/stacks-chow/pinned-original-source.tex line 14734).\par\medskip
\noindent\textit{Allowed standard prerequisites (editorial).}\begin{itemize}
\item Cycles and their presentation as sums of integral closed subschemes. Same chapter, Definition labelled definition-cycles, cited at proof lines 14752-14754. Provides the decomposition alpha = sum n\_i [X\_i] with g : coprod X\_i -> X proper.
\item Flat pullback, Gysin maps and proper pushforward as bivariant classes. Same chapter, Lemma labelled lemma-push-proper-bivariant, cited at proof lines 14776-14778, with the flat pullback and Gysin classes used at proof lines 14758-14768. Builds xi\_i in the two cases and shows that g\_* o xi is a bivariant class.
\item Testing bivariant classes on integral closed subschemes. Same chapter, Lemma labelled lemma-bivariant-zero, cited at proof lines 14792-14794. Reduces the inverse identity to the case of an integral X'.
\item The two shapes of an integral scheme over a one dimensional base. Same chapter, the construction of ord and the cycle theory for the situation, invoked at proof lines 14796-14802. Gives [X'] = (p')\textasciicircum{}*[S]\_1 when p' is flat of relative dimension dim\_delta(X') - 1, and [X'] = (p')\textasciicircum{}*(s -> S)\textasciicircum{}*[S]\_1 otherwise.
\end{itemize}\medskip
\ifdefined\hypertarget\hypertarget{sourcepointer-1}{}\fi
\noindent\textit{Source reference.} Original definition 02QR: sources/stacks-chow/pinned-original-source.tex, source lines 1535--1535. Same-chapter pointer into the author environment definition at source line 1535; invoked at proof line 14757 of the counted proof. It is retained only as a navigation pointer; every load-bearing step is inside the delivered span.\par\medskip
\ifdefined\hypertarget\hypertarget{sourcepointer-2}{}\fi
\noindent\textit{Source reference.} Original lemma 0EPK: sources/stacks-chow/pinned-original-source.tex, source lines 6167--6167. Same-chapter pointer into the author environment lemma at source line 6167; invoked at proof line 14781 of the counted proof. It is retained only as a navigation pointer; every load-bearing step is inside the delivered span.\par\medskip
\ifdefined\hypertarget\hypertarget{sourcepointer-3}{}\fi
\noindent\textit{Source reference.} Original lemma 0FC4: sources/stacks-chow/pinned-original-source.tex, source lines 14658--14658. Same-chapter pointer into the author environment lemma at source line 14658; invoked at proof line 14784 of the counted proof. It is retained only as a navigation pointer; every load-bearing step is inside the delivered span.\par\medskip
\ifdefined\hypertarget\hypertarget{sourcepointer-4}{}\fi
\noindent\textit{Source reference.} Original lemma 02UC: sources/stacks-chow/pinned-original-source.tex, source lines 6585--6585. Same-chapter pointer into the author environment lemma at source line 6585; invoked at proof line 14797 of the counted proof. It is retained only as a navigation pointer; every load-bearing step is inside the delivered span.\par\medskip
\begin{lemma}
\label{lemma-chow-cohomology-towards-base-dim-1}
Let $(S, \delta)$ be as above. Let $X$ be a scheme locally of finite type
over $S$. Then we have a canonical identification
$$
A^p(X \to S) = \CH_{1 - p}(X)
$$
for all $p \in \mathbf{Z}$.
\end{lemma}

\begin{proof}
Consider the element $[S]_1 \in \CH_1(S)$. We get a map
$A^p(X \to S) \to \CH_{1 - p}(X)$ by sending $c$ to $c \cap [S]_1$.

\medskip\noindent
Conversely, suppose we have $\alpha \in \CH_{1 - p}(X)$.
Then we can define $c_\alpha \in A^p(X \to S)$ as
follows: given $X' \to S$ and $\alpha' \in \CH_n(X')$
we let
$$
c_\alpha \cap \alpha' = \alpha \times \alpha'
$$
in $\CH_{n - p}(X \times_S X')$. To show that this is a bivariant
class we write $\alpha = \sum_{i \in I} n_i[X_i]$ as in
Definition \ref{definition-cycles}. In particular the morphism
$$
g : \coprod\nolimits_{i \in I} X_i \longrightarrow X
$$
is proper. Pick $i \in I$. If $X_i$ dominates an irreducible component
of $S$, then the structure morphism $p_i : X_i \to S$ is flat and we have
$\xi_i = p_i^* \in A^p(X_i \to S)$. On the other hand, if $p_i$ factors
as $p'_i : X_i \to s_i$ followed by the inclusion $s_i \to S$
of a closed point, then we have
$\xi_i = (p'_i)^* \circ c_i \in A^p(X_i \to S)$
where $c_i \in A^1(s_i \to S)$ is the Gysin homomorphism and
$(p'_i)^*$ is flat pullback. Observe that
$$
A^p(\coprod\nolimits_{i \in I} X_i \to S) =
\prod\nolimits_{i \in I} A^p(X_i \to S)
$$
Thus we have
$$
\xi = \sum n_i \xi_i \in A^p(\coprod\nolimits_{i \in I} X_i \to S)
$$
Finally, since $g$ is proper we have a bivariant class
$$
g_* \circ \xi \in A^p(X \to S)
$$
by Lemma \ref{lemma-push-proper-bivariant}. The reader easily
verifies that $c_\alpha$ is equal to this class
(please compare with the proof of
Lemma \ref{lemma-exterior-product-well-defined-dim-1})
and hence is itself a bivariant class.

\medskip\noindent
To finish the proof we have to show that the two constructions
are mutually inverse. Since $c_\alpha \cap [S]_1 = \alpha$
this is clear for one of the two directions. For the other, let
$c \in A^p(X \to S)$ and set $\alpha = c \cap [S]_1$.
It suffices to prove that
$$
c \cap [X'] = c_\alpha \cap [X']
$$
when $X'$ is an integral scheme locally of finite type over $S$,
see Lemma \ref{lemma-bivariant-zero}. However, either $p' : X' \to S$
is flat of relative dimension $\dim_\delta(X') - 1$ and hence
$[X'] = (p')^*[S]_1$ or $X' \to S$ factors as $X' \to s \to S$
and hence $[X'] = (p')^*(s \to S)^*[S]_1$. Thus the fact that the
bivariant classes $c$ and $c_\alpha$ agree on $[S]_1$
implies they agree when capped against $[X']$ (since bivariant classes
commute with flat pullback and Gysin maps) and the proof is complete.
\end{proof}

\end{document}
